\documentclass[a4paper,12pt,french]{report}

% =======================
% IMPORTATION DES PACKAGES
% =======================

\usepackage[T1]{fontenc}
\usepackage[utf8]{inputenc}
\usepackage[margin=2cm]{geometry}
\usepackage{theorem}
\usepackage{fourier}
\usepackage{amssymb, amsmath, amsfonts, dsfont}
\usepackage{babel}

\renewcommand{\familydefault}{\sfdefault}

\pagestyle{headings}

\ifpdf
\usepackage[pdftex=true, hyperindex=true, colorlinks=false]{hyperref}
\else
\usepackage[pdftex=true, hyperindex=true, colorlinks=true]{hyperref}
\fi

\theoremstyle{break}
\theoremheaderfont{\scshape}
\theorembodyfont{\upshape}

\newtheorem{defin}{Définition}[section]
\newtheorem{theo}[defin]{Théorème}
\newtheorem{prop}[defin]{Proposition}
\newtheorem{coro}[defin]{Corollaire}
\newtheorem{lemme}[defin]{Lemme}
\newtheorem{rema}[defin]{Remarque}


\title{MÉSURE ET INTÉGRATION}
\date{\today}
\author{}

\begin{document}
\maketitle

\textsc{Précision} \\

Ceci est juste des notes pour m'aider à comprendre la théorie de la mesure et intégration tout en apprenant latex.
Mais si cela peut vous aider, vous pouvez le consulter et le télécharger.

\tableofcontents % table de matiere

% --- THEORIE DES ENSEMBLES ---

\chapter{THÉORIE DES ENSEMBLES}

\section{Théorie des ensembles}
\begin{defin}
Soit $X$ un ensemble non vide et $(A_{i})_{i \in \mathbb{N}}.$ 
      $$\bigcup_{i \in \mathbb{N}} A_{i} := \{ x \in X | \mbox{ il existe } i_{0} \mbox{ tel que } x \in A_{i_{0}}\}.$$
      $$\bigcap_{i \in \mathbb{N}} A_{i} := \{ x \in X | \mbox{ pour tout } i_{0} \mbox{ tel que } x \in A_{i_{0}}\}.$$
      $$A \backslash B := \{ x \in A \mbox{ et } x \not\in B\}.$$
      $$X\backslash A := \{ x \not\in A\}.$$
      $$ A \triangle B := (A \backslash  B) \cap (B \backslash A).$$ 
\end{defin}

\begin{prop}[Loi de Morgan]
     $$X\backslash \left(\bigcup_{i \in \mathbb{N}} A_{i}\right) = \bigcap_{i \in \mathbb{N}} (X \backslash A_{i}).$$
     $$X\backslash \left(\bigcap_{i \in \mathbb{N}} A_{i}\right) = \bigcup_{i \in \mathbb{N}} (X \backslash A_{i}).$$
\end{prop}

\begin{prop}[Distributivité]
     $$A \cap (B \cup C) = (A \cap B) \cup (A \cap C).$$
     $$A \cup (B \cap C) = (A \cup B) \cap (A \cup C).$$
\end{prop}

\begin{defin}
       $$\mathds{1}_{A}(x) = \left\{ \begin{array}{c} 1 \mbox{ si } x \in A \\ 0 \mbox{ si } x \not\in A\end{array}\right.$$
\end{defin}

\begin{prop}
       \begin{enumerate}
              \item $\mathds{1}_{X\backslash A} = 1 - \mathds{1}_{A}$,
              \item $\mathds{1}_{A \cap B} = \mathds{1}_{A} \mathds{1}_{B}$,
              \item $\mathds{1}_{A \cup B} = \mathds{1}_{A} + \mathds{1}_{B} - \mathds{1}_{A \cap B}$,
              \item $\mathds{1}_{A \backslash B} = \mathds{1}_{A} - \mathds{1}_{A \cap B}$,
              \item $\mathds{1}_{A \triangle B} = \mathds{1}_{A} + \mathds{1}_{B} \mbox{ mod } 2$.
       \end{enumerate}
\end{prop}

\section{Ensemble ordonnée}
\begin{defin}[Relation d'ordre]
       Soit X un ensemble non vide. On dit que  $\preccurlyeq$ est une relation d'ordre s'il verifie :
       \begin{itemize}
              \item[i)] Reflexivité : Soit $x \in X$, $x \preccurlyeq x$,
              \item[ii)] Anti-symetrie : Soit $x, y \in X$, $x \preccurlyeq y$ et $y \preccurlyeq x$, alors $x=y$,
              \item[iii)] Transitivité : Soit $x, y, z \in X$, $x \preccurlyeq y$ et $y \preccurlyeq z$, on a $x \preccurlyeq z$.   
       \end{itemize}
       On appelle $(X, \preccurlyeq)$ un ensemble ordonnée.
\end{defin}

\begin{defin}[Relation d'ordre total et partielle]
       On dit que l'ensemble ordonnée $(X, \preccurlyeq)$ est total si tout élement de X est comparable c'est à dire pour $x, y \in X$, $x \preccurlyeq y$ ou bien $y \preccurlyeq x$.
       Si c'est le contraire, on dit qu'elle est partielle.
\end{defin}

% --- TRIBU, CLAN ET CLASSE MONOTONE ---

\chapter{TRIBUS, CLAN ET CLASSE MONOTONE}
\section{Terminologie}
\begin{defin}[Clan]
      Soit X un ensemble non vide. On dit que $\mathcal{A}$ est un clan sur X si elle verifie :
      \begin{itemize}
             \item $\emptyset \in \mathcal{A}$,
             \item $\mathcal{A}$ est stable par passage au complémentaire, c'est à dire pour $A \in \mathcal{A}$, on a $X \backslash A \in \mathcal{A}$,
             \item $\mathcal{A}$ est stable par réunion, c'est à dire pour $A, B \in \mathcal{A}$, on a $A \cup B \in \mathcal{A}$.
      \end{itemize}
\end{defin}

\begin{prop}
	Soit X un ensemble non vide. L.a.s.s.e
	\begin{enumerate}
		\item $\mathcal{A}$ est un clan sur X.
		\item
			\begin{itemize}
				\item[i)] $\mathcal{A}$ est stable par intersection (resp. réunion) finie sur X,
				\item[ii)] $\mathcal{A}$ est stable par différence ensembliste.
			\end{itemize}
	\end{enumerate}
\end{prop}

\begin{defin}[Tribus]
      Soit X un ensemble non vide. On dit que $\mathcal{A}$ est une tribu sur X si elle vérifie :
      \begin{itemize}
             \item $\emptyset \in \mathcal{A}$,
             \item $\mathcal{A}$ est stable par passage au complémentaire, c'est à dire pour $A \in \mathcal{A}$, on a $X \backslash A \in \mathcal{A}$,
             \item $\mathcal{A}$ est stable par réunion dénombrable, c'est à dire pour $(A_{n})_{n \in \mathbb{N}} \in \mathcal{A}$, on a $\displaystyle \bigcup_{n \in \mathcal{N}} A_{n} \in \mathcal{A}$.
      \end{itemize}
\end{defin}

\begin{prop}
        Soit X un ensemble non vide. L.a.s.s.e
        \begin{enumerate}
                \item $\mathcal{A}$ est un tribus sur X.
                \item
                        \begin{itemize}
                                \item[i)] $\mathcal{A}$ est stable par intersection (resp. réunion) dénombrable sur X,
                                \item[ii)] $\mathcal{A}$ est stable par différence ensembliste.
                        \end{itemize}
        \end{enumerate}
\end{prop}

\begin{defin}[Classe monotone]
       On dit que $\mathcal{A}$ est une classe monotone si elle vérifie : 
       \begin{itemize}
              \item $(A_{n})_{n \in \mathbb{N}} \nearrow A$, alors $\displaystyle\bigcup_{n \in \mathbb{N}} A_{n} \in \mathcal{A}$
              \item $(A_{n})_{n \in \mathbb{N}} \searrow A$, alors $\displaystyle\bigcap_{n \in \mathbb{N}} A_{n} \in \mathcal{A}$
       \end{itemize}
\end{defin}

\begin{theo}
	Soit $\mathcal{A}$ un clan sur X. Alors L.a.s.s.e
	\begin{itemize}
		\item[i)] $\mathcal{A}$ est un tribus sur X,
		\item[ii)] $\mathcal{A}$ est un classe monotone sur X.
	\end{itemize}
\end{theo}

\begin{defin}[Préclan]
        On dit que $\mathcal{A}$ est un préclan sur $X$ si:
        \begin{itemize}
              \item[i)] $\emptyset \in \mathcal{A}$,
              \item[ii)] $\mathcal{A}$ est stable par intersection finie,
              \item[iii)] Pour tous $A,B \in \mathcal{A}$, $B \backslash A$ est une réunion disjoints dans $\mathcal{A}$, c'est-à-dire il existe $I \subset \mathbb{N}$ et une suite $(A_{i})_{i \in \mathbb{N}} \subset \mathcal{A}$, on a $B \backslash A = \displaystyle \bigsqcup_{i \in \mathbb{N}} A_{i}$
        \end{itemize}
\end{defin}

\section{Structure engendré}
\begin{lemme}
         Tout intersection d'un clan (resp. tribus, resp. classe monotone) est un clan (resp. tribus, resp. classe monotone).
\end{lemme}

\chapter{MESURABILITÉ}
\section{Espace mesurable}
\begin{defin}
       Soit X un ensemble non vide. On dit que $(X,\mathcal{A})$ est un espace mesurable si $\mathcal{A}$ est un tribus.
\end{defin}

\begin{prop}[Fonction mesurable]
       \begin{itemize}
              \item Soit $(X, \mathcal{A})$ et $(Y, \mathcal{B})$ deux espace mesurable et $f : X \mapsto Y$ une fonction.\\ On dit que $f$ est $\mathcal{A}-\mathcal{B}$-mesurable ou mesurable si pour tout $B \in \mathcal{B}, \; f^{-1}(B) \in \mathcal{A}$.
              \item Si $X$ et $Y$ sont muni des tribus borelien respectif \textit{B}($X$) et \textit{B}($Y$). Lorsque $f$ est $\mathcal{A}-\mathcal{B}$-mesurable, on dit qu'elle est \textit{B}($X$)-\textit{B}{$Y$}-borelienne ou borelienne.
       \end{itemize}
\end{prop}

\begin{rema}
       \begin{itemize}
              \item Soient $(X, \mathcal{A})$ et $(Y, \mathcal{B})$ deux espace mesurable $f : X \mapsto Y$ une application. Si f est mesurable on le note $\mathcal{L}(X,Y)$ ou $\mathcal{L}(\mathcal{A}, \mathcal{B})$. 
              \item Toute fonction continue est mesurable.
       \end{itemize}
\end{rema}

\begin{theo}[Tribus initiale]
	La tribu $f^{-1}(B)$ pour $B \in \mathcal{B}$ est la plus petit tribu qui rende f $\mathcal{A}-\mathcal{B}$-mesurable.
\end{theo}

\begin{theo}[Principe de recollement]
	Soit $(X, \mathcal{A})$ et $(Y, \mathcal{B})$ deux espace mesurable, $(A_{n})_{n \in \mathbb{N}} \subseteq \mathcal{A}$ une partition denombrable de X et $f : X \mapsto Y$ une application. Si $f_{| A_{i}}$ est $\mathcal{A}_{A_{i}}-\mathcal{B}$-mesurable, alors $f$ est $\mathcal{A}-\mathcal{B}$-mesurable.
\end{theo}

% --- CHAPITRE MESURE ---

\chapter{MESURE}
\section{Géneralité sur les mesures}
On tavaille sur un espace mesurable $(X, \mathcal{A})$.

\begin{defin}[Application additive vs $\sigma$-additive]
\begin{itemize}
       \item Additive: $\mu (A \cup B) := \mu (A) + \mu (B)$, pour tous $A, B \in \mathcal{A}$ disjoints.
       \item $\sigma$-additivé: Pour toute suite d'ensemble disjoints $(A_{n})_{n \in \mathbb{N}} \subseteq \mathcal{A}$, on a: $$\mu (\bigsqcup_{n \in \mathbb{N}} A_{n}) := \sum_{n \in \mathbb{N}} \mu (A_{n})$$.
\end{itemize}
\end{defin}

\begin{defin}[Définition d'une mesure]
Une application $\mu$ sur $\mathcal{A}$ est une mesure (positive) si elle est à valeurs dans $[0, +\infty]$, verifiant $\mu (\emptyset)$ et $\sigma$-additivé.
\end{defin}

\begin{prop}[Propriété fondamentale de la mesure]
\begin{itemize}
       \item Croissance: Si $A \subseteq B$, alors $\mu (A) \leq \mu (B)$.
       \item $\sigma$-additivé: Pour toute suite quelconque $a_{n}$, $$\mu (\bigcup_{n \in \mathbb{N}} A_{n}) \leq \sum_{n \in \mathbb{N}} \mu (A_{n}).$$
       \item Continuité par modification monotone:
       \begin{itemize}
              \item Si $A_{n} \nearrow \mathcal{A}$, alors $\displaystyle \lim_{n \rightarrow \infty} \mu (A_{n}) = \mu (\bigcup_{n \in \mathbb{N}} A_{n})$, 
              \item Si $A_{n} \searrow \mathcal{A}$, alors $\displaystyle \lim_{n \rightarrow \infty} \mu (A_{n}) = \mu (\bigcap_{n \in \mathbb{N}} A_{n})$.
       \end{itemize}
\end{itemize}
\end{prop}

\begin{prop}[Mesure image par une application]
Si $f : X \rightarrow Y$ est une application $\mathcal{A}-\mathcal{B}$-mesurable, l'application $\nu (B) = \mu(f^{-1}(B))$ définit une mesure sur $\mathcal(B)$ appelé mesure image.
\end{prop}

\begin{prop}[Mesure finie et $\sigma$-finie]
\begin{itemize}
       \item Finie: Si $\mu (X) \le \infty$,
       \item $\sigma$-finie: S'il existe une suite $(A_{n})_{n \in \mathbb{N}} \subseteq \mathcal{A}$ telle que $X = \bigcup_{n \in \mathbb{N}} A_{n}$ avec $\mu (A_{n} \le \infty)$ pour tout $n$.
\end{itemize}
\end{prop}

\begin{defin}[Continuité absolue et etrangeté]
\begin{itemize}
       \item $\mu$ est absolument continue par rapport a $\nu \; (\mu \ll \nu)$ si $\nu(A) = 0$, on a $\mu (A) = 0$.
       \item $\mu$ et $\nu$ sont etrangere ($\mu \perp \nu$) s'il existe une partition de $X = A \cup (X \backslash A)$ telle que $\nu(A)=0$ et $\mu(X \backslash A) = 0$.
\end{itemize}
\end{defin}

\section{Classes ensembles negligeable et completion}
\begin{defin}[Ensemble $\mu$-negligeable]
Un ensemble $N \subseteq X$ est negligeable s'il est contenue dans un élement $A \in \mathcal{A}$ telle que $\mu (A) = 0$.
\end{defin}

\begin{defin}[Espace mesure complet]
Un espace $(X, \mathcal{A}, \mu)$ est dit complet si tout ensemble $\mu$-negligeable appartient à la tribu $\mathcal{A}$.
\end{defin}

\begin{defin}[Tribu completée $\mathcal{A}_{\mu}$]
C'est la plus petit tribus contenant $\mathcal{A}$ et tous les ensembles $\mu$-negligeables.\\ Un ensemble $A \in \mathcal{A}_{\mu}$ s'écrit sous la forme $A = B \cup N$ ou $B \in \mathcal{A}$ et $N$ est negligeable.
\end{defin}

\begin{theo}[Mesure completée]
La mesure $\mu$ s'etend de manière unique à $\mathcal{A}_{\mu}$ en posant $\mu (B \cup N) = \mu (B)$.
\end{theo}

\begin{prop}[Propriété presque partout ($\mu-p.p$)]
Une propriété est vrai $\mu$-p.p si elle est vraie partout sauf sur un ensemble $\mu$-negligeable.
\end{prop}

\section{Construction de mesure par Carathéodory}
\begin{prop}[Mesure exterieur]
       Une application $\mu^{*} : P(X) \rightarrow [0, \infty]$ telle que:
       \begin{enumerate}
              \item $\mu^{*}(\emptyset) = 0$,
              \item $A \subset B \Rightarrow \mu^{*}(A) \leq \mu^{*}(B)$,
              \item $\displaystyle \mu^{*}(\bigcup_{n \in \mathbb{N}} A_{n}) \leq \sum_{n \in \mathbb{N}} \mu^{*}(A_{n})$.
       \end{enumerate}
\end{prop}

\begin{prop}[Ensemble $\mu^{*}$-mesurable]
       Un ensemble $A \subset X$ est dit $\mu^{*}$-mesurable si pour tout ensemble de test $B \subset X$, $\mu^{*} = \mu^{*}(B \cap A) + \mu^{*}(B \cap A^{c})$.
\end{prop}

\begin{theo}[Théorème de Carathéodory]
       L'ensemble des parties $\mu^{*}$-mesurable forme une tribu et la restriction $\mu^{*}$ à cette tribu est une vraie mesure.
\end{theo}

\section{Mesure borelienne et mesure de Lebesgue}
C'est la mise en pratique sur la droite réelle $\mathbb{R}$.

\begin{defin}[Mesure de Borel-Stieltjes]
       Associe à une fonction $\Phi : \mathbb{R} \rightarrow \mathbb{R} \nearrow$ et continue à gauche.
       Elle est definie sur les intervalles par : $\mu_{\Phi}([a, b[) = \Phi(b) - \Phi(a)$
\end{defin}

\begin{defin}[Mesure de Borel]
       C'est le cas particulier où $\Phi (x) = x$ (fonction identité). Elle est définie sur la tribu Borelienne $\mathcal{B}(\mathbb{R})$.
\end{defin}

\begin{defin}[Mesure de Lebesgue $\lambda$]
       C'est la completée de la mesure de Borel sur la tribu de Lebesgue $\mathcal{M}_{\mathbb{R}}$
\end{defin}

\begin{prop}[Proprieté incontournable de la mesure de Lebesgue]
       \begin{itemize}
              \item La mesure d'un singleton est nulle $\lambda(\{ x \})=0$,
              \item Pour tout intervalle borné de borne $a$ et $b$, sa mesure et sa longueur : $$\lambda([a,b[) = \lambda(]a,b[) = \lambda([a,b]) = b-a .$$
              \item Invariance par translation: pour tout borelien $A$ et tout $x \in \mathbb{R}, \; \lambda(A + x) = \lambda(A)$.
              \item Invariance par symetrie: $\lambda(-A) = \lambda(A)$.
       \end{itemize}
\end{prop}

% --- CHAPITRE INTEGRALE ---

\chapter{INTÉGRALE}
\section{Les définition fondamentale}
\begin{defin}[Fonction etagée (ou simple)]
       Soit $(X, \mathcal{T}, \mu)$ un espace mesuré. Une fonction $f : X \rightarrow [0, \infty[$ est dite etagee si elle est mesurable et ne prend qu'un nombre finie de valeurs positive distincts $\{ \alpha_{1}, \alpha_{2}, \cdots, \alpha_{n}\}$.\\ Elle s'ecrit sous la forme $ \displaystyle f(x) = \sum_{i=1}^{n} \alpha_{i} \mathds{1}_{A_{i}}(x).$ où $A_{i} = s^{-1}(\{ \alpha_{i}\})$ sont des ensemble mesurables disjoints.
\end{defin}

\begin{defin}[Intégrale d'une fonction etagée positive]
       L'intégrale de $f$ par rapport à la mesure $\mu$ est définie par $ \displaystyle \int_{X} f d\mu = \sum_{i=1}^{n} \alpha_{i} \mu(A_{i})$ avec la convention $0 \times \infty = 0$.
\end{defin}

\begin{defin}[Intégrale d'une fonction mesurable positive]
       Soit $f ; X \rightarrow [0, \infty]$ une fonction mesurable positive de Lebesgue. Son intégrale est définie par la borne superieure des integrales des fonction etagée qui la dominent par en dessous : $$\int_{X} f d\mu = \mbox{ sup } \left\{ \int_{X} s d\mu \; | \; 0 < s < f \mbox{ s est etagée} \right\}$$.
\end{defin}

\begin{defin}[Fonction intégrale de signe quelconque]
       Soit $f : X \rightarrow \overline{\mathbb{R}} \mbox{ ou } \mathbb{C}$ une fonction mesurable. On définit ses parties positive et négative : $f^{+} = max(f, 0)$ et $f^{-} = max(-f, 0)$. On dit que $f$ est intégrable au sens de Lebesgue (on note $f \in \mathcal{L}^{1}(\mu)$) si $$\int_{X} |f|d\mu < +\infty \Leftrightarrow \int_{X} f^{+} d\mu < +\infty \mbox{ et } \int_{X} f^{-}d\mu < +\infty.$$
       Si $f \in \mathcal{L}^{1}(\mu)$, son intégrable est : $ \displaystyle \int_{X} fd\mu = \int_{X} f^{+}d\mu - \int_{X} f^{-} d\mu$.
\end{defin}

\section{Théorème fondamenataux de passage a la limite}
\begin{theo}[convergence monotone]
       Soit $(f_{n})_{n \in \mathbb{N}}$ une suite de fonction mesurable positive de $X$ dans $[0, \infty]$. Si la suite est $\nearrow$ p.p, c'est a dire $0 \leq f_{1}(x) \leq \cdots \leq f_{n}(x) \mu$-p.p. Alors la fonction limite $f(x) = \lim_{n \rightarrow +\infty} f_{n}(x)$ est mesurable et $$\lim_{n \rightarrow +\infty} \int_{X} f_{n} d\mu = \int_{X}\lim_{n \rightarrow +\infty} f_{n} d\mu = \int_{X} f d\mu.$$ 
\end{theo}

\begin{lemme}[Fatou] 
       Soit $(f_{n})_{n \in \mathbb{N}}$ une suite de fonction mesurable positive sur $X$ dans $[0, +\infty]$. Alors $$\int_{X} \lim_{n \rightarrow +\infty} inf f_{n} d\mu \leq \lim_{n \rightarrow +\infty} inf \int_{X} f_{n} d\mu.$$
\end{lemme}

\begin{theo}[Convergence dominé ou théorème de Lebesgue]
       Soit $(f_{n})_{n \in \mathbb{N}}$ une suite de fonction mesurable (de signe quelconque) à valeurs dans $\mathbb{R}$ ou $\mathbb{C}$.\\ Supposons que :
       \begin{enumerate}
              \item La suite $(f_{n})$ converge $\mu$-p.p vers une fonction $f$.
              \item (Hypothese de domination)\\ Il existe une fonction $g \in \mathcal{L}^{1}(\mu)$ positive et intgrable telle que $|f_{n}(x)| \leq g(x) \forall n \in \mathbb{N} \mu$-p.p.
       \end{enumerate}
       Alors $f$ est intégrable et on peut intervertir la limite: $$\lim_{n  \rightarrow +\infty} \int_{X} f_{n} d\mu = \int_{X} f d\mu.$$
\end{theo}

\section{Proposition importantes}
\begin{prop}[Linéarité et positivité]
       L'application $f \rightarrow \int_{X} f d\mu$ est linéaire sur $\mathcal{L}^{1}(\mu)$ et croissante (si $f \leq g \mu$-p.p, alors $\int f \leq \int g$).
\end{prop}

\begin{prop}[Negligeable et intégration]
       \begin{itemize}
              \item Si $f$ est une fonction mesurable positive et $\int_{X} f d\mu = 0$, alors $f=0 \mu$-p.p.
              \item Si $f \in \mathcal{L}^{1}(\mu)$, alors $|f(x)| < +\infty \mu$-p.p.
       \end{itemize}
\end{prop}

\begin{prop}[Intégration sur les sous ensembles]
       Soit $A \in \mathcal{T}$ un ensemble mesurable.\\ L'intégrale de $f$ sur $A$ est définie par $\displaystyle \int_{A} f d\mu = \int_{X} f \mathds{1}_{A} d\mu$.\\ Si $\mu(A)=0$, alors pour toute fonction mesurable $f$, $\displaystyle \int_{A} f d\mu = 0$.
\end{prop}

\section{Corollaire cruciaux pour la pratique}
\begin{coro}[Intégrale des serie et fonction positive]
       Soit $(u_{n})_{n \in \mathbb{N}}$ une suite de fonction mesurable positives. Alors $\displaystyle \int_{X} \sum_{n=0}^{+\infty} u_{n} (x) d\mu = \sum_{n=0}^{+\infty} \int_{X} u_{n} (x) d\mu$
\end{coro}

\begin{coro}[Intégration des series de fonction quelconque]
       Soit $(u_{n})_{n \in \mathbb{N}}$ une suite de fonction mesurable a valeurs réelle ou complexe. Si la serie des intégrales des valeurs absolues converge: $$\sum_{n=0}^{+\infty} \int_{X} |u_{n}|d\mu < +\infty.$$
       Alors la serie $\displaystyle \sum_{n=0}^{+\infty}$ converge $\mu$-p.p, sa somme est intégrable et $$\int_{X} \sum_{n=0}^{+\infty} u_{n} d\mu = \sum_{n=0}^{+\infty} \int_{X} u_{n} d\mu.$$
\end{coro}

\section{Théorème à paramètre (continuité et derivation)}
\begin{theo}[Continuité sous les signes]
       Soit $I$ un intervalle de $\mathbb{R}$. Si
       \begin{itemize}
              \item[i)] Pour $\mu$-presque tout $x \in X$, la fonction $t \mapsto f(t,x)$ est continue sur $I$,
              \item[ii)] Pour tout $t \in I$, la fonction $x \mapsto f(t,x)$ est mesurable,
              \item[iii)] (Domination local) \\Pour tout compact $K \subset I, \; \exists \; g_{K} \in \mathcal{L}^{1}(\mu)$ positive telle que $|f(t,x)| \leq g_{K}(x) \; \forall t \in K \mu$-p.p.
       \end{itemize}
       Alors $\displaystyle F(t) = \int_{X} f(t,x) d\mu(x)$ est continue sur $I$.
\end{theo}

\begin{theo}[Dérivation sous le signe $\int$]
       Soit $I$ un intervalle de $\mathbb{R}$. Si
       \begin{itemize}
              \item[i)] $\forall \; t \in I, \; x \mapsto f(t,x)$ est intégrable,
              \item[ii)] Pour $\mu$-presque tout $x \in X, \; x \mapsto f(t,x)$ est derivable sur $I$ (de dérivé partielle $\frac{\partial}{\partial t} f(t,x)$),
              \item[iii)] (Domination de al derivé) \\ $\exists \; g \in \mathcal{L}^{1}(\mu)$ positive telle que $\frac{\partial}{\partial t} f(t,x) \leq g(x) \; \forall t \in I, \; \mu$-p.p. 
       \end{itemize}
       Alors $F$ est derivable sur $I$, $\frac{\partial}{\partial t} f(t,\cdot)$ est intégrable pour tout $t$ et \begin{align*}  F^{'}(t)  & =  \frac{d}{dt} \int_{X} f(t,x) d\mu(x) \\ & =  \int_{X} \frac{\partial}{\partial t} f(t,x)d\mu(x)\end{align*}
\end{theo}
\end{document}
